This chapter will show the maths...
\section{Plane Waves}
Light is electromagnetic radiation. From Maxwell's equations it can be shown that the electric field, $\mathbf{E}$, component of electromagnetic radiation satisfies the wave equation
\begin{equation}
\nabla^2 \mathbf{E} = \frac{1}{c^2}\frac{\partial^2}{\partial t^2}\mathbf{E}.
\label{eq:wave_equation}
\end{equation}
A monochromatic, coherent, plane wave solution for $\mathbf{E}$ satisfies equation (\ref{eq:wave_equation}):
\begin{equation}
\mathbf{E} = \mathbf{\hat{n}}E \exp\big( i (\omega t - \mathbf{k}.\mathbf{r} + \varphi) \big),
\label{eq:plane_wave_solution}
\end{equation}
where $\mathbf{\hat{u}}$ indicates the polarisation, $E$ is the complex field amplitude, $\omega$ is the angular frequency, $\mathbf(k)$ is the wave vector and $\varphi$ is the absolute phase. The physical amplitude of the electric field oscillation is obtained from the $\Re(E)$. For optical systems is useful to define an optical axis. When the light is travelling along the optical axis, the $ \mathbf{k}.\mathbf{r}$ becomes $kz$. $k$ is known as the wave number and is related to the wavelength of the light,$\lambda$ by 
\begin{equation}
k = \frac{2\pi}{\lambda}.
\end{equation}
For plane waves there is no dependence of $\mathbf{E}$ on the $x$ and $y$ coordinates. The absolute phase $\varphi$ becomes relevant in calculations involving interfering beams. By inserting equation (\ref{eq:plane_wave_solution}) into equation (\ref{eq:wave_equation}), we obtain the dispersion relationship 
\begin{equation}
\omega = ck.
\end{equation}
The plane wave solution can be used to analyse many optical phenomena, such as amplitude modulation, phase modulation and interference.  
\section{Gaussian Modes}
Light in an optical set-up is normally generated by a laser. Lasers produce almost monochromatic, coherent, light so in this sense they are quite similar to a plane wave, however lasers also have geometric properties. A laser produces light with an approximately circular beam profile which is brightest at its centre. The spread of light power in a laser beam follows a Gaussian distribution, and the width of the beam profile is characterised by the radius at which $\sim86\%$ of the light power is contained. 
\paragraph{}
As the beam propagates along its optical axis, the beam size changes. A Gaussian beam has a point where it is smallest, this is known as the waist of the beam. A Gaussian beam can be completely characterised by the size of its waist and location along the waist along the optical axis. Whenever a Gaussian mode interacts with a simple optical component such as a lens or a mirror, the current Gaussian mode is transformed into a new Gaussian mode.
\paragraph{}
The following sections describe Gaussian modes and how they interact with real optical layouts with imperfections. The lowest order Gaussian mode will be derived, and the properties of this mode will be explained. The concept of the complex Gaussian beam parameter will be introduced as a way of encapsulating all the information needed to determine the properties of a Gaussian mode. Astigmatic Gaussian beams will be defined in terms of two complex Gaussian beam parameters. The first order Hermite-Gaussian (HG) mode will be shown to model the effect of misalignment, and then it will be described how all laser beams can made up of a set of HG modes. Finally, the way that a Gaussian mode traverses an optical set-up will be expressed in terms of the ABCD matrix formalism.
\section{The Paraxial Wave Equation}
\paragraph{Short intro on what this is}
\paragraph{}
To derive the paraxial wave equation, we apply the separation of variables to equation (\ref{eq:wave_equation}):
\begin{equation}
\mathbf{E}(\mathbf{r},t) = E(\mathbf{r})T(t).
\end{equation}
Next we reform equation \ref{eq:wave_equation} in terms of $E$ and $T$
\begin{align}
\frac{\nabla^2 E}{E} = \frac{1}{c^2T}\frac{\partial^2 T}{\partial T^2}.
\label{eq:sep_of_var}
\end{align}
By noticing that either side of the equation (\ref{eq:sep_of_var}) depends on $\mathbf{r}$ or $t$, the left hand side of this equation must equal a constant, and the Helmholtz equation can be obtained:
\begin{equation}
\frac{\nabla ^2 E}{E} = -k^2.
\label{eq:helmholtz_equation}
\end{equation}
Rewriting $E(\mathbf{r})$ in terms of two functions,
\begin{equation}
E(\mathbf{r}) = u(\mathbf{r})\exp\big(-ikz\big),
\label{eq:generic_mode}
\end{equation} 
and substituting this into equation (\ref{eq:helmholtz_equation}) yields
\begin{equation}
\Bigg(\left( \frac{\partial^2}{\partial x^2} + \frac{\partial^2}{\partial y^2} \right) u(\mathbf{r}) - 2ik \frac{\partial u}{\partial z}  + \frac{\partial^2 u}{\partial z^2} \Bigg)\exp\big(ikz\big) = 0.
\end{equation}
The paraxial approximation means the $\frac{\partial^2 u}{\partial z^2}$ is neglected as it is much smaller than the other terms. Applying this approximation produces the paraxial Helmholtz equation
\begin{equation}
\left(\frac{\partial^2}{\partial x^2} + \frac{\partial^2}{\partial y^2} \right) u(\mathbf{r}) - 2ik \frac{\partial u}{\partial z} = 0.
\label{eq:paraxial_helmholtz}
\end{equation} 
The physical interpretation of $u(\mathbf{r})$ is a function that describes the spatial distribution of the complex amplitude of the light field. \textcolor{red}{reword this so it actually makes sense.}  
\section{The fundamental Gaussian mode}
The simplest geometry of a laser beam can be described by the fundamental Gaussian mode. This is a mode which has a circular intensity profile that follows a Gaussian distribution, and the surface of constant phase is curved. Substituting the simplest solution to (\ref{eq:paraxial_helmholtz}) into equation (\ref{eq:generic_mode}) produces the fundamental Gaussian mode \textcolor{red}{find out what the amplitude constant is}
\begin{equation}
E(\mathbf{r}) = E_0\sqrt{\frac{2}{\pi}} \frac{1}{w(z)} \exp\big(i\phi(z)\big)\exp\left( -ik\frac{x^2 + y^2}{2R_C(z)} - \frac{x^2 + y^2}{w^2(z)}\right)\exp\big(-ikz\big).
\label{eq:gaussian_mode}
\end{equation}
This is characterised by a function $w(z)$, which gives the size of the beam at position $z$ the optical axis, by $\phi(z)$, which gives the Gouy Phase of the beam at $z$ along the optical axis, and by $R_c(z)$, which gives radius of curvature of the beam at $z$. The minimum value of $w(z)$ is known as the beam waist, or $w_0$.
\paragraph{}
It is useful to express $w(z)$,$\phi(z)$ and $R_c(z)$ in terms of the Rayleigh range $z_{\rm{R}}$. The Rayleigh range characterises how fast the beam will spread out. A beam with a short Rayleigh length will expand quicker than a beam with a longer Rayleigh length. For example, in order for the beam to not be significantly clipped by the aperture of the test masses, the Rayleigh length of the Gaussian beam in the arm cavities at LIGO has to be on the order of kilometres. If the waist of a Gaussian beam at $z = 0 $, the position along the $z$ axis by which the radius has increased by a factor of $\sqrt(2)$ is known as the Rayleigh range $z_{\rm{R}}$, and is given by
\begin{equation}
z_{\rm{R}} = \frac{\pi w_0^2}{\lambda}.
\end{equation} 
\paragraph{}
The radius of the beam as it propagates along the optical axis can be written as
\begin{equation}
	w(z) = w_0\sqrt{1 + \left(\frac{z}{z_{\rm{R}}}\right)^2}.
	\label{eq:beam_size}
\end{equation} 
The Rayleigh range can be used to approximately split beam into two regions: the near field for $z<z_{\rm{R}}$ and the far field for $z>z_{\rm{R}}$. In the far field equation (\ref{eq:beam_size}) becomes linear and the beam size can be approximated as 
\begin{equation}
w(z) \approx w_0\frac{z}{z_{\rm{R}}}.
\end{equation}
\paragraph{}
As a Gaussian beam travels propagates along the optical axis, it acquires an additional phase shift compared to a plane wave. This is known as the Gouy phase shift, $\phi(z)$, and it given by
\begin{equation}
	\phi(z) = \tan^{-1}\frac{z}{z_{\rm{R}}}.
\end{equation}
At either side of it waist, the fundamental Gaussian mode will be shifted $\pm pi/2$ relative to a plane wave. An alternative way of putting this is the fundamental Gaussian mode will gain a $\pi$ change as it goes through from $-\infty$ to $\infty$ compared to a plane wave. Most of this phase change is gained within $\pm z_{\rm{R}}$ of its waist. The Gouy phase results in the the phase velocity at the waist being greater for a Gaussian beam than it is for plane waves.
\paragraph{}
The radius of curvature, $R_c(z)$, in terms of the Rayleigh range is
\begin{equation}
R = z + \frac{z_{\rm{R}}^2}{z}.
\end{equation}
For $z \ll z_{\rm{R}}$ the radius of curvature can be approximated as infinite. As the beam approaches its Rayleigh length the radius of curvature decreases until it reaches its minimum at $z = z_{\rm{R}}$. After this the radius of curvature begin to flatten out again, and for $z \gg z_{\rm{R}}$ the radius of curvature can be approximated as $z$. The defocus, $S(z)$, of a beam is a more practical way of describing a beam as it is easier to compute the effects of a beam interacting with a focussing element of an optical layout in terms of the defocus than it is with the radius of curvature. The defocus is defined as the inverse of its radius of curvature
\begin{equation}
	S(z) = \frac{1}{R(z)}
\end{equation}
\paragraph{}
The fundamental Gaussian mode can be expressed in terms of a complex beam parameter, $q(z)$. This parameter is useful for calculating how a gaussian beam is transformed by an optical system. From this parameter, all the physical quantities related to a Gaussian beam can be computed. The complex beam parameter is defined as
\begin{equation}
q(z) = z + iz_{\rm{R}}.
\label{eq:eq_q}
\end{equation}
The inverse of equation \ref{eq:eq_q} can be written in terms of the curvature and radius of the beam:
\begin{equation}
\frac{1}{q(z)} = \frac{1}{R(z)} - i\frac{\lambda}{\pi w^2(z)}.
\end{equation} 
Equation (\ref{eq:gaussian_mode}) can be re written in terms of this the complex beam parameter,
\begin{equation}
E(\mathbf{r}) = E_0\sqrt{\frac{2}{\pi}} \frac{q_0}{w_0q(z)} \exp\big(i\phi(z)\big)\exp\left( -ik\frac{x^2 + y^2}{2q(z)}\right)\exp\big(-ikz\big).
\end{equation}
The radius of the beam, gouy phase, and the defocus at a point $z$ along the optical axis can be calculated from $q(z)$ using
\begin{align}
w(z) & =  \sqrt{\frac{\lambda}{-\pi\Im\left(\frac{1}{q(z)}\right)}},\\
\phi(z) &= \tan^{-1}\left(\frac{\Re\big(q(z)\big)}{\Im\big(q(z)\big)}\right),\\
S(z) &= \Re\left( \frac{1}{q(z)} \right).
\end{align}
\paragraph{}

\paragraph{Mode matching integral}
\section{Astigmatism}
\paragraph{Wordy definition of simple astigmatism}
\section{Gaussian Beams Traversing optical Layouts}
\paragraph{talk about the concept of tracing the beam parameter through the optical layout}
\paragraph{ABCD matrices for optical elements}
\section{Higher Order Modes and optical imperfections}
\paragraph{wordy description of what a HOM is}
\paragraph{simple example of misalignment}
\paragraph{Coupling or Scattering to higher order modes}
\paragraph{mathematical description of HG modes}
\paragraph{More general HOM expansion}
\paragraph{point out Gouy phase}
